\documentclass[a4paper,9pt]{extarticle}

%=========================================================
% PACKAGES
%=========================================================

\usepackage[
a4paper,
top=1.8cm,
bottom=1.5cm,
footskip=7mm,
left=1cm,
right=1cm,
headheight=14pt,
headsep=5mm
]{geometry}
\usepackage[T1]{fontenc}
\usepackage[utf8]{inputenc}
\usepackage[french]{babel}

\usepackage[sfdefault]{FiraSans}
\renewcommand{\familydefault}{\sfdefault}

\usepackage{multicol}
\setlength{\columnsep}{0.7cm}

\usepackage{amsmath,amssymb}

\usepackage{titlesec}
\usepackage{enumitem}

\usepackage{tcolorbox}
\tcbuselibrary{skins,breakable}

\usepackage{xcolor}

\usepackage{graphicx}
\usepackage{float}

\usepackage{tikz}

\usepackage{fancyhdr}
\usepackage[hidelinks]{hyperref}


%=========================================================
% PAGE SETTINGS
%=========================================================

\setlength{\parindent}{0pt}
\setlength{\parskip}{1.5pt}

\setlength{\abovedisplayskip}{4pt}
\setlength{\belowdisplayskip}{4pt}

\setlist[itemize]{
leftmargin=*,
itemsep=1pt,
topsep=1pt
}


%=========================================================
% COLORS
%=========================================================

\definecolor{Blue}{HTML}{E3F0FF}
\definecolor{Green}{HTML}{E3F7EA}
\definecolor{Pink}{HTML}{FFE5ED}
\definecolor{Orange}{HTML}{FFF0DC}
\definecolor{Purple}{HTML}{EEE7FF}
\definecolor{Yellow}{HTML}{FFFADB}
\definecolor{Teal}{HTML}{E0F5F3}


\definecolor{BlueTitle}{HTML}{285B8F}
\definecolor{GreenTitle}{HTML}{287052}
\definecolor{PinkTitle}{HTML}{A84063}
\definecolor{OrangeTitle}{HTML}{A96500}
\definecolor{PurpleTitle}{HTML}{6245A5}
\definecolor{TealTitle}{HTML}{1F6F68}


%=========================================================
% HEADER
%=========================================================

\pagestyle{fancy}

\fancyhf{}

\fancyhead[L]{\small\textbf{Modèles stochastiques}}
\fancyhead[C]{\small Ergodisme, densité spectrale, systèmes LTI}
\fancyhead[R]{\small Week 05 --- 05-11/10/2026}

\fancyfoot[C]{\small\href{https://github.com/AlphaXT38G/notes_ba5}{github.com/AlphaXT38G/notes\_ba5}}
\fancyfoot[R]{\small \thepage}

\renewcommand{\headrulewidth}{0.4pt}



%=========================================================
% TITLES
%=========================================================

\titleformat{\section}
{\large\bfseries\color{BlueTitle}}
{}{0pt}{}

\titlespacing{\section}
{0pt}{6pt}{3pt}


% separator

\newcommand{\separator}{
\vspace{2pt}
\hrule
\vspace{3pt}
}


%=========================================================
% BOXES
%=========================================================

% Boxes are kept whole (unbreakable): inside multicols, a box split by tcolorbox
% leaves an empty column behind it.

% DEFINITION - blue
\newtcolorbox{definition}{
enhanced,
unbreakable,
colback=Blue,
colframe=BlueTitle,
boxrule=0.7pt,
arc=0mm,
title=\textbf{Definition},
fonttitle=\bfseries
}

% IMPORTANT - pink
\newtcolorbox{important}{
enhanced,
unbreakable,
colback=Pink,
colframe=PinkTitle,
boxrule=0.7pt,
arc=0mm,
title=\textbf{Key idea},
fonttitle=\bfseries
}

% FORMULAS - green
\newtcolorbox{formula}{
enhanced,
unbreakable,
colback=Green,
colframe=GreenTitle,
boxrule=0.7pt,
arc=0mm,
title=\textbf{Key formulas},
fonttitle=\bfseries
}

% WARNING - orange
\newtcolorbox{warning}{
enhanced,
unbreakable,
colback=Orange,
colframe=OrangeTitle,
boxrule=0.7pt,
arc=0mm,
title=\textbf{Warning},
fonttitle=\bfseries
}

% EXAMPLE - purple
\newtcolorbox{example}{
enhanced,
unbreakable,
colback=Purple,
colframe=PurpleTitle,
boxrule=0.7pt,
arc=0mm,
title=\textbf{Example},
fonttitle=\bfseries
}

% SUMMARY - orange
\newtcolorbox{summary}{
enhanced,
unbreakable,
colback=Yellow,
colframe=OrangeTitle,
boxrule=0.7pt,
arc=0mm,
title=\textbf{Summary},
fonttitle=\bfseries
}



% SIMPLE - teal: the same idea again, in the plainest words, with an everyday picture
\newtcolorbox{simple}{
enhanced,
unbreakable,
colback=Teal,
colframe=TealTitle,
boxrule=0.7pt,
arc=0mm,
title=\textbf{Expliqué simplement},
fonttitle=\bfseries
}



%=========================================================
% FIGURE BOX
%=========================================================

\newtcolorbox{figurebox}{
enhanced,
colback=white,
colframe=gray,
boxrule=0.5pt,
arc=0mm,
center
}



%=========================================================
% DOCUMENT
%=========================================================

\begin{document}

%\footnotesize   --> To reduce the size of the text

\thispagestyle{fancy}



\begin{multicols}{2}


%=========================================================
\section{Où on en est}


\begin{summary}

\begin{itemize}

\item Mardi : module 3, ergodisme.

\item Vendredi : module 3, densité spectrale, réponse d'un système LTI à des signaux
aléatoires.

\item Exercices M3 : 9 (mardi) ; 10, 11 (vendredi).

\end{itemize}

\end{summary}


\begin{important}

\textbf{Conseils du Moodle.}

\begin{itemize}

\item \textbf{Exercice 10} (difficile, prépare le 11) : une transformation linéaire $Y = AX$ de
gaussiennes multivariées donne des gaussiennes multivariées. Au point b, il faut trouver la
matrice $A$ qui \textbf{décorrèle} les deux gaussiennes : $Y$ est alors formé de 2 gaussiennes
indépendantes.

\item \textbf{Exercice 11} : passer de $\operatorname{Arctan}(\dots)$ à
$\operatorname{Arcsin}(\dots)$ demande des manipulations trigonométriques fastidieuses, on peut
garder la réponse en $\operatorname{Arctan}$.

\end{itemize}

\end{important}


%=========================================================
\section{Ergodisme par rapport à la moyenne}


\begin{definition}

$X(t)$ WSS. On estime la moyenne d'ensemble $\mu_X = E[X(t)]$ par la \textbf{moyenne
temporelle} d'une seule réalisation :

\[
\langle X(t)\rangle_T = \frac1T\int_0^TX(t)\,dt
\]

$X$ est \textbf{ergodique par rapport à sa moyenne} si, quand $T\to\infty$,

\[
E[\langle X(t)\rangle_T]\to\mu_X \qquad \mathrm{VAR}[\langle X(t)\rangle_T] \to 0
\]

(alors $\langle X(t)\rangle_T$ converge en moyenne quadratique vers $\mu_X$).

\end{definition}


\begin{formula}

La première condition est toujours vraie (estimateur \textbf{non biaisé}) : $E[\langle
X\rangle_T] = \frac1T\int_0^T\mu_X\,dt = \mu_X$. Pour la seconde :

\[
\mathrm{VAR}[\langle X\rangle_T] = \frac1{T^2}\int_0^T\!\!\int_0^TC_X(t-s)\,dt\,ds
\]
\[
= \frac2T\int_0^TC_X(\tau)\Big(1 - \frac\tau T\Big)d\tau
\]

\end{formula}


\begin{formula}

\textbf{Théorème 1} (nécessaire et suffisant) : $X$ WSS est ergodique par rapport à sa moyenne
ssi

\[
\lim_{T\to\infty}\frac1T\int_0^TC_X(\tau)\Big(1 - \frac\tau T\Big)d\tau = 0
\]

\textbf{Corollaire 1} (suffisant) : si $\lim_{\tau\to\infty}C_X(\tau) = 0$, $X$ est ergodique
par rapport à sa moyenne.

\end{formula}


\begin{important}

\textbf{Preuve du corollaire.} Pour $\varepsilon > 0$, il existe $T_0$ tel que $|C_X(\tau)| <
\varepsilon$ pour $\tau > T_0$. On coupe l'intégrale en $T_0$ et on borne $|C_X| \le C_X(0)$
avant $T_0$ :

\[
\frac1T\Big|\int_0^TC_X(\tau)\big(1-\tfrac\tau T\big)d\tau\Big| \le C_X(0)\frac{T_0}T +
\varepsilon\Big(1 - \frac{T_0}T\Big)
\]

qui tend vers $\varepsilon$, et ce pour tout $\varepsilon > 0$.

\end{important}


\begin{simple}

Deux façons de faire une moyenne : demander à \textbf{1000 personnes} leur valeur au même
instant (moyenne d'ensemble), ou suivre \textbf{une seule personne} très longtemps (moyenne
temporelle). Ergodique : les deux donnent la même chose, une seule longue mesure suffit. Ça
marche si le processus « oublie » son passé : quand deux instants très éloignés ne sont plus
liés ($C_X\to0$), une longue trajectoire explore tout.

\end{simple}


%=========================================================
\section{Autres formes d'ergodisme et exemples}


\begin{definition}

\textbf{Par rapport à la variance} :

\[
E[\langle X^2\rangle_T - \mu_X^2]\to\sigma_X^2 \qquad \mathrm{VAR}[\langle X^2\rangle_T -
\mu_X^2]\to0
\]

\textbf{Par rapport à l'auto-corrélation}, avec $\langle X(t+\tau)X(t)\rangle_T = \frac1T
\int_0^TX(t+\tau)X(t)\,dt$ :

\[
E[\langle X(t+\tau)X(t)\rangle_T] \to R_X(\tau)
\]
\[
\mathrm{VAR}[\langle X(t+\tau)X(t)\rangle_T]\to0
\]

On se ramène au cas de la moyenne avec $Z(t,\tau) = X(t+\tau)X(t)$.

\vspace{3pt}

\textbf{En distribution} (plus fort, pour un SSS) : les estimées temporelles des fonctions de
répartition convergent en moyenne quadratique vers elles.

\end{definition}


\begin{example}

\begin{itemize}

\item \textbf{Sinusoïde à phase uniforme} : ergodique par rapport à sa moyenne si $f_0 \ne 0$
(chaque trajectoire est une sinusoïde centrée en 0, de moyenne temporelle nulle), et à sa
variance (exercice 5). Si $f_0 = 0$ : $X(t,\varphi_0) = a\sin\varphi_0$ est constante, sa
moyenne temporelle vaut $a\sin\varphi_0 \ne 0$ : \textbf{pas} ergodique.

\item \textbf{Phase et amplitude aléatoires} $X(t) = A\sin(2\pi f_0t + \Phi)$, $A\in\{a_1,a_2\}$
(Bernoulli) : ergodique par rapport à sa moyenne, mais \textbf{pas} à sa variance (certaines
trajectoires ont l'amplitude $a_1$, d'autres $a_2$).

\item \textbf{Séquence binaire aléatoire} : ergodique par rapport à sa moyenne par le
corollaire ($C_X(\tau) = 0$ pour $|\tau| \ge T$).

\end{itemize}

\end{example}


\begin{simple}

Le contre-exemple à retenir : si chaque trajectoire « choisit » au départ quelque chose qu'elle
garde pour toujours (une amplitude, une constante), suivre une seule trajectoire ne montre
qu'\textbf{un} choix. On ne peut pas deviner les autres : pas ergodique.

\end{simple}


%=========================================================
\section{Densité spectrale de puissance}


\begin{formula}

$X(t)$ WSS. \textbf{Densité spectrale de puissance} (relations de \textbf{Wiener-Khintchine}) :

\[
S_X(f) = \int_{-\infty}^{\infty}R_X(\tau)\,e^{-j2\pi f\tau}\,d\tau
\]
\[
R_X(\tau) = \int_{-\infty}^{\infty}S_X(f)\,e^{j2\pi f\tau}\,df
\]

\textbf{Puissance moyenne} :

\[
E[X^2(t)] = R_X(0) = \int_{-\infty}^\infty S_X(f)\,df
\]

\end{formula}


\begin{warning}

Bien définie si $\int|R_X(\tau)|\,d\tau < \infty$. Faux si $\mu_X \ne 0$, car $R_X(\tau) =
C_X(\tau) + \mu_X^2$. Alors soit on ajoute un Dirac, $S_X(f) = \mathcal F(C_X)(f) +
\mu_X^2\delta(f)$, soit (de préférence) on travaille avec le processus centré $X - \mu_X$.
Certains processus centrés ne la vérifient quand même pas : ceux à dépendances à long terme
(fin du cours).

\end{warning}


\begin{formula}

\textbf{Propriétés} ($X$ réel) :

\begin{itemize}

\item \textbf{P1.} $S_X(f)$ est \textbf{réelle et paire} (car $R_X$ est paire).

\item \textbf{P2.} $S_X(f) \ge 0$ pour tout $f$.

\item \textbf{P3.} $h(f) = S_X(f)/R_X(0)$ est une densité de probabilité.

\end{itemize}

\textbf{Densité spectrale mutuelle} : $S_{XY}(f) = \mathcal F(R_{XY})(f)$, en général
complexe.

\end{formula}


\begin{simple}

$S_X(f)$ dit \textbf{combien de puissance} le processus met à chaque fréquence, comme
l'égaliseur d'une chaîne hi-fi : beaucoup de graves, peu d'aigus... Additionner sur toutes les
fréquences donne la puissance totale. Et on l'obtient simplement en prenant Fourier de $R_X$.

\end{simple}

\end{multicols}


\end{document}
